Um zu zeigen, dass die Zylinderfläche eine reelle Mannigfaltigkeit ist, betrachten wir zunächst die Konstruktion des Zylinders durch die Verklebung der Kanten eines Quadrats. Sei das Quadrat in der Ebene durch die Punkte \((0,0)\), \((1,0)\), \((1,1)\) und \((0,1)\) gegeben. Die Verklebung der vertikalen Kanten \((0,y)\) und \((1,y)\) für \(y \in [0,1]\) erfolgt durch die Identifikation \((0,y) \sim (1,y)\). Die horizontale Verklebung ist nicht notwendig, um einen Zylinder zu erhalten, da sie bereits durch die Identifikation der vertikalen Kanten impliziert ist.

Ein Atlas für den Zylinder kann durch zwei Karten gegeben werden. Die erste Karte \((U_1, \varphi_1)\) sei definiert auf der offenen Menge \(U_1 = (0,1) \times (0,1)\) mit der Kartenabbildung \(\varphi_1: U_1 \to \mathbb{R}^2\), \(\varphi_1(x,y) = (x,y)\). Die zweite Karte \((U_2, \varphi_2)\) deckt die Umgebung der verklebten Kanten ab und sei definiert durch \(U_2 = (-\epsilon, 1+\epsilon) \times (0,1)\) für ein kleines \(\epsilon > 0\), mit der Kartenabbildung \(\varphi_2: U_2 \to \mathbb{R}^2\), \(\varphi_2(x,y) = (\cos(2\pi x), \sin(2\pi x), y)\).

Die Übergangsfunktionen zwischen diesen Karten sind glatt, da die Verklebung der Kanten keine Singularitäten einführt. Insbesondere ist die Übergangsfunktion \(\varphi_2 \circ \varphi_1^{-1}\) auf dem Schnitt \(U_1 \cap U_2\) gegeben durch die Identität, da \(\varphi_1\) und \(\varphi_2\) in diesem Bereich übereinstimmen.

Um zu prüfen, ob dieser Atlas eine komplexe Struktur definiert, müssen die Übergangsfunktionen holomorph sein. Da die Übergangsfunktionen in diesem Fall lediglich die Identität sind, sind sie trivialerweise holomorph. Somit definiert der gegebene Atlas auch eine komplexe Struktur auf der Zylinderfläche.

Zusammenfassend ist die Zylinderfläche eine reelle Mannigfaltigkeit, und der angegebene Atlas definiert auch eine komplexe Struktur.